\title{Large Language Models Can Automatically Engineer Features for Few-Shot Tabular Learning}

\begin{abstract}
Large Language Models (LLMs) have demonstrated outstanding proficiency in addressing complex reasoning tasks that they have not explicitly encountered before, presenting a transformative opportunity for tabular learning—an area fundamental to numerous practical applications. This study introduces a new in-context learning paradigm, termed \textsf{FeatLLM}, where LLMs assume the role of automated feature engineers, generating modified datasets that maximize performance in downstream tabular prediction problems. The novel features crafted by the LLM are subsequently leveraged by a straightforward predictive model (e.g., linear regression), supporting superior few-shot learning capabilities. Notably, \textsf{FeatLLM} requires only the output of this predictive model with the features derived from LLMs at inference, obviating the need for continuous querying of the LLM for each test instance as required by alternative LLM-based methods. Additionally, \textsf{FeatLLM} demands nothing more than API-level integration with LLMs and deftly avoids the restrictions imposed by limited prompt sizes. Experiments spanning a diverse array of tabular benchmarks validate the efficacy of \textsf{FeatLLM}, with the generated rule sets considerably surpassing competing techniques such as TabLLM and STUNT—achieving an average performance gain of 10\%.
\end{abstract}

\section{Introduction}
Large Language Models (LLMs) have lately established themselves as powerful engines capable of responding effectively to entirely new tasks, all without relying on specific fine-tuning for those tasks~\cite{creswell2022selection,imani2023mathprompter,taylor2022galactica}. By being trained on vast textual datasets, these models internalize extensive background knowledge, allowing them to serve as practical surrogates for specialized human expertise in numerous domains~\cite{Genomes2021,singhal2023large,wilcox2003role}. This unique attribute has paved the way for automation advancements in areas as varied as dialogue understanding~\cite{finch2023leveraging} and the automated repair of software~\cite{fan2023automated}. One particularly impactful property of LLMs is their ability, when supplied with concise task statements and a handful of examples, to interpret the contextual cues of a problem and to focus on the most salient features and instances~\cite{brown2020language,zhao2023survey}. These behaviors strongly indicate that LLMs can provide expert-level feature engineering by means of data analysis. 

Leveraging the inherent knowledge and advanced reasoning skills of LLMs is increasingly influential in the landscape of tabular data analysis. To illustrate, contextual knowledge—such as recognizing age as a risk factor for certain illnesses—can improve predictive accuracy in medical domains. Recent research trends have encompassed designing prompts that describe tasks and features in plain language, serializing table data, and presenting these inputs to LLMs~\cite{dinh2022lift,wang2023anypredict}. Such approaches are commonly followed by in-context learning protocols~\cite{nam2023semi} or modest fine-tuning techniques that require minimal parameter updates~\cite{dinh2022lift,hegselmann2023tabllm}. Harnessing LLMs' ingrained expertise has thus been shown to outperform standard machine learning methods in data-scarce (low-shot) scenarios~\cite{hegselmann2023tabllm}.

Present LLM-based methods for learning from tabular data exhibit several shortcomings. For end-to-end prediction tasks, at least one inference through the LLM must be performed per individual instance, resulting in substantial computational overhead. Achieving strong predictive performance typically necessitates fine-tuning of the LLM, rendering these methods incompatible with many LLMs that do not offer accessible training environments or APIs for adaptation—a significant obstacle given that many state-of-the-art LLMs allow only restricted usage via inference APIs~\cite{anil2023palm,openai2023gpt4}. Moreover, most existing approaches struggle with prompt length limitations~\cite{liu2023lost}: when tabular datasets feature numerous columns, the textual representation often becomes unwieldy, either making the solution impractical or leading to degraded effectiveness. Such limitations introduce fundamental obstacles to deploying LLM-driven tabular learning solutions in practical contexts.

In contrast, our proposed methodology reframes the use of LLMs in tabular tasks: rather than using LLMs exclusively for end-to-end output prediction, we employ them for feature construction, leveraging their extensive prior knowledge more efficiently. We present an innovative in-context learning paradigm, \textsf{FeatLLM}, designed to extract and formalize the “criteria” utilized by LLMs in their decision-making processes. For example, when tasked with predicting whether a patient has a specific disease, the LLM can deduce and articulate rules linking certain feature values to the disease prediction. These articulated criteria enable us to generate novel features that can substitute the original input features in subsequent machine learning workflows. This strategic shift repositions the LLM as a feature generator, enabling the downstream model to be simpler—once informative features are available, inference no longer relies on a complex model, thus enhancing computational efficiency. The feature extraction process capitalizes on the in-context learning strengths of LLMs, requiring only carefully crafted prompts with demonstration examples, without the need for actual model retraining. Additionally, leveraging ensemble-inspired techniques such as feature bagging~\cite{breiman2001random} enables us to address issues arising from oversized prompts.

\textsf{FeatLLM} organizes the prompt for creating novel features into two distinct stages: (1) comprehending the underlying problem and deducing connections between features and target labels; (2) leveraging insights from the prior stage alongside a few examples to infer discriminative rules that can separate different classes. This systematic reasoning encourages LLMs to disregard irrelevant features during the process of rule formulation. The formulated rules are executed over the dataset (translated as programmatic instructions), converting the data into binary features that denote the fulfillment of each rule by individual samples. Subsequently, a simple model is trained using these derived features to approximate class probabilities. To enhance robustness, an ensemble technique is adopted where feature generation is repeated with feature bagging. By finely tuning the selection of feature subsets and sample counts for bagging, the method addresses issues related to overly lengthy prompts. As \textsf{FeatLLM} does not require additional model training, it incurs minimal inference overhead, making LLM-based approaches applicable for diverse practical tasks.

Our experiments on 13 tabular datasets within the low-shot learning paradigm demonstrate that \textsf{FeatLLM} achieves consistently strong and resilient outcomes. The results indicate that LLMs effectively harness both existing knowledge and information obtained from the dataset to construct high-quality features. Across multiple scenarios, our method surpasses current few-shot learning benchmarks. The implementation is made publicly available via an anonymized GitHub repository at~\texttt{https://github.com/Sungwon-Han/FeatLLM}.

\section{Related Work}

\subsection{Few-Shot Learning with Tabular Data}
Few-shot learning has progressed to a stage where it enables the extraction of transferable representations from data, significantly reducing the cost associated with manual labeling~\cite{chen2018closer}. Originally developed for image-based tasks, the applicability of few-shot learning has since expanded to encompass other modalities such as text, audio, and, more recently, tabular data~\cite{majumder2022few,nam2023stunt,schick2021exploiting}. Relying on a limited set of labeled samples introduces the possibility of capturing misleading associations~\cite{bartlett2021deep}. In response, current methodologies often integrate external sources of information or infuse domain-oriented prior knowledge to embed the right inductive biases into the model training procedure. Strategies include semi-supervised learning by leveraging extensive collections of unlabeled tabular data combined with thoughtful augmentation~\cite{ucar2021subtab,yoon2020vime}, or improving model initialization through broadly-applicable unsupervised pre-training~\cite{somepalli2021saint}. Unsupervised training objectives typically use techniques such as masking or intentionally corrupting sections of the data to enforce a reconstruction objective~\cite{arik2021tabnet,majmundar2022met}, or utilize data augmentation to define contrastive learning losses~\cite{bahri2022scarf,chen2020simple}. The intrinsic heterogeneity present in tabular datasets, however, limits the general applicability of these augmentation approaches, and they have not demonstrated substantial success in few-shot circumstances~\cite{nam2023stunt}.

To address these limitations, recent research proposes that models first be exposed to a diverse spectrum of tabular datasets—beyond those similar to the intended target task—after which transfer learning is employed~\cite{zhang2023generative,levin2022transfer,zhu2023xtab}. As an illustration, the TransTab approach~\cite{wang2022transtab} makes use of transformer architectures trained to recognize columnar semantics using column names in addition to raw tabular values. This allows extracted knowledge about columns and tables to support zero-shot or few-shot generalization for new tasks. Separately, TabPFN~\cite{hollmann2022tabpfn} synthesizes varied distributions to generate artificial data for the pre-training of Bayesian neural networks. Once pre-trained, the model infers on new tasks by jointly inputting both labeled and unlabeled samples—requiring no fine-tuning—yielding performance superior to classical tree-based models, especially when labeled examples are scarce.

\subsection{Language-Interfaced Tabular Learning}
An alternative path for introducing prior knowledge in tabular learning relies on the integration of large language models (LLMs)~\cite{anil2023palm,openai2023gpt4}. Due to their extensive training on heterogeneous textual data, LLMs have exhibited the capacity to generalize effectively to novel challenges across many application areas~\cite{brown2020language}. To transfer the encoded knowledge of LLMs to tabular domains, recent techniques focus on reformatting tabular datasets as textual input sequences, which are then embedded within LLM prompts~\cite{dinh2022lift,hegselmann2023tabllm,wang2023anypredict}. Prompt construction may include not just the raw tabular entries, but also metadata such as feature descriptions and task statements, enabling the model to reason about feature-task relations during inference. There is an ongoing trend toward further harnessing this knowledge using either parameter-efficient training paradigms such as LoRA~\cite{hu2021lora} or IA3~\cite{liu2022few}, or alternatively through in-context learning, which makes use of annotated examples provided in the input prompt to enable prediction without updating model weights~\cite{dinh2022lift,hegselmann2023tabllm,nam2023semi,slack2023tablet}.

Although previous approaches have proven useful in advancing few-shot tabular learning, their reliance on running inference through the LLM for every prediction introduces obstacles, particularly for deployment in industrial settings. This research takes a different direction by foregoing the typical black-box pipeline where each sample's output is generated end-to-end using an LLM. Rather, the focus is shifted towards having the LLM deduce explicit conditions or rule sets for each class. \textsf{FeatLLM} interprets these rules into executable code to construct supplementary features, thereby facilitating predictions independently from the LLM. This approach utilizes in-context learning to extract class-specific rules based on prior knowledge and a limited number of labeled instances.

\section{Methods}
\textsf{FeatLLM} is architected to derive ``rules''—the criteria defining each class's answers—from the limited input examples, as shown in Figure~\ref{fig:main_model}. The LLM processes the input samples to generate the rules, which are then converted to code used to formulate new binary features. These features signal the presence or absence of rule satisfaction for each sample. The resulting binary features are combined through a dot product with non-negative, trainable weights to predict the probability of membership in each class. Training this model involves adjusting the weights to reflect how much each feature contributes to the class probability (Section~\ref{Sec:training}). This workflow is reiterated multiple times with bagging and the outputs are merged using an ensemble. The specific problem setting and the granular steps of this methodology are discussed in detail below.

\vspace*{-0.12in}
\paragraph{Problem formulation.} Consider a tabular dataset consisting of $N$ annotated instances, $\mathcal{D} = \{(\mathbf{x}^i, \mathbf{y}^i)\}_{i=1}^{N}$. In this dataset, each entry has $d$ input attributes, meaning each $\mathbf{x}^i$ is represented as a vector of length $d$. The attributes are identified by their natural language names, collected as $F = \{f_j\}_{j=1}^{d}$, and each comes with a separate definition. In the context of classification, each label $\mathbf{y}^i$ corresponds to a member of a predetermined set of classes. In $k$-shot learning scenarios, a subset of $k$ ($k < N$) labeled examples are sampled randomly for model training.

\subsection{Prompt Design for Extracting Rules}
\label{Sec:prompt_design}
In order to help the LLM derive rules using more robust reasoning strategies, we structure the problem-solving approach to reflect methods that an expert would use when analyzing tabular data. The prompt itself consists of three main sections (see Fig.~\ref{fig:prompt_for_rule} and Appendix~\ref{sec:example_rule} for details):

\vspace*{-0.12in}
\paragraph{Basic information description.} This section delivers the crucial information required to tackle the task (refer to orange highlights in Fig.~\ref{fig:prompt_for_rule}). It incorporates the task specification with label data, descriptions of each feature, and illustrative demonstrations. The task specification is framed as a direct question (e.g., ``Does this patient's myocardial infarction complication data suggest chronic heart failure? Yes or no?"). Additional examples are available in Appendix~\ref{sec:task_desc}. Feature descriptions clarify the value types and may also include definitions. For categorical features, representative categories are given as examples. Demonstration examples entail converting several training instances into textual form together with their correct labels. The serialization utilizes formats established in prior literature~\cite{dinh2022lift,hegselmann2023tabllm}:

\begin{align}
&\text{Serialize}(\mathbf{x}^i,\mathbf{y}^i, F) = \nonumber \\ 
&``\ f_1\ \text{is}\ \mathbf{x}^i_1.\ ... \ f_d\  \text{is}\  \mathbf{x}^i_d.\ \ \text{Answer:}\  \mathbf{y}^i\ ”
\end{align}

\vspace*{-0.12in}
\paragraph{Reasoning instruction.} Our goal is to improve the reasoning abilities of LLMs by incorporating reasoning instructions into the prompt (as highlighted in blue in Fig.~\ref{fig:prompt_for_rule}). The reasoning instruction opens with a sentence reminiscent of the chain-of-thought prompting technique~\cite{wei2022chain}. Subsequently, we structure the LLM's inference workflow in two phases. Initially, the LLM is tasked with analyzing the causal relationships or underlying tendencies between given features and the task description. In this phase, the model is directed to leverage general knowledge or commonsense reasoning, deliberately avoiding reference to example demonstrations, which enables the extraction of background information relevant to the task. The second phase involves applying the learned general knowledge and example demonstrations from the first phase to construct rules corresponding to each class. This sequential reasoning process steers the model away from relying on irrelevant feature correlations and helps prioritize features that are most consequential. To strike a balance between underfitting and overly verbose prompts, we establish a constant rule count—specifically, 10 rules\footnote{Section~\ref{sec:analysis} contains further discussion regarding the selection of rule quantity.} for each class. Furthermore, we instruct the LLM to consider feature types during rule generation by referring to the basic information description, thereby minimizing the risk of rule-type mismatches.

\vspace*{-0.12in}
\paragraph{Response instruction.} For improved parsing and practical utility of the generated rules, we delineate explicit response instructions for the LLM (see the yellow annotations in Fig.~\ref{fig:prompt_for_rule}). These instructions specify both the overall format required for responses by answer class (i.e., response format) and the particular structure rules should conform to (i.e., rule format). The guidelines empower the LLM to choose suitable formats consistent with the feature type. In the absence of more restrictive directives, the LLM may compose compound rules using logical connectors such as AND or OR. Illustrative rule outcomes are found in Appendix~\ref{sec:outcome-rule}.

\subsection{Inferring Class Likelihood via Rules}
\label{Sec:training}

\paragraph{Feature extraction from parsed rules.} In this stage, the rules produced previously are employed to derive novel binary features. Each feature corresponds to a particular class, flagging whether a sample meets the criteria stipulated by rules for that class (denoted by a value of ‘0’ or ‘1’). Such features contribute to estimating the likelihood that a given sample belongs to each class. Because the LLM formulates rules using natural language, an automated parsing mechanism is required to reliably convert textual rules into features. Although the rule-generation step restricts the response format and rule structure to streamline parsing, LLM outputs frequently exhibit unpredictable syntactic variations~\cite{kovavc2023large}. As an example, expressions like ``$>$” or ``$<$” may be replaced with their word equivalents such as ``greater than" or ``less than," complicating direct parsing.

To overcome difficulties posed by noisy rule texts, rather than developing intricate parsing algorithms, we delegate parsing to the LLM itself. Given its proficiency in grasping semantic content regardless of syntactic variability, and its ability to output simple code instructions (given its extensive code-oriented training corpus), the LLM can efficiently generate functional parsing code. Our prompts to the LLM consist of a function identifier, explanations for inputs and outputs, as well as the derived rules; these prompts are then submitted to the LLM for code generation. Figures~\ref{fig:prompt_for_parsing} and ~\ref{sec:example_parsing}-\ref{sec:example_parse_outcome} in Appendix illustrate example prompts and the resultant code snippets. The code returned by the LLM is subsequently executed through Python's \texttt{exec()} command to perform data transformation.

\vspace*{-0.12in}
\paragraph{Assessing class probabilities.} The binary features distilled from rules for each class permit a basic approach for quantifying class membership probabilities: summing the number of rules a sample satisfies for each class (that is, the aggregate value of binary features by class). Yet, since the significance of individual rules and features can vary, it becomes important to determine their relative importance from training data. We propose to do this by employing standard machine learning (ML) models tailored to each class's binary feature vector. For illustration, a linear model without bias may be used. Let the binary feature vector for sample $\mathbf{x}^i$ corresponding to class $k$ be $\mathbf{z}_k^i \in \{0, 1\}^R$, where $R$ denotes the count of rules associated with a given class and $c$ the total number of classes. The probability $\mathbf{p}^i$ for each class can be derived as follows:

\begin{align}
\text{logit}_k^i &= \max(\mathbf{w}_k, 0) \cdot \mathbf{z}_k^i, \nonumber \\
\mathbf{p}^i &= \text{Softmax}({[\text{logit}_{1}^i, ..., \text{logit}_{c}^i]}). \label{eq:infer_prob}
\end{align}

In this context, $\mathbf{w}_k$ denotes the tunable parameters within the linear model, representing the contribution of each input feature. By constraining these weights to lie within a positive domain, it is possible to guarantee that features generated by the LLM do not decrease the probability of a class during training. The logits are mapped to class probabilities using the softmax transformation. The parameters $\mathbf{w}_k$ are learned through the minimization of the cross-entropy loss, using a few-shot dataset.

\vspace*{-0.12in}
\paragraph{Ensembling with bagging.} To form the final ensemble, the procedure is iterated multiple times, each yielding a distinct inference model. For $T$ independent runs, the resulting weight vectors $\{\mathbf{w}[t]\}_{t=1}^T$ are used to produce predictions by calculating the mean of the class probabilities $\mathbf{p}^i$ derived for each $\mathbf{w}[t]$, as defined by Eq.~\ref{eq:infer_prob}.

To foster diversity in rule generation for the ensemble, the LLM inference is performed at a high temperature, and the sequence of few-shot exemplars is randomly permuted; this is motivated by findings that the order of demonstrations can affect LLM outputs~\cite{lu2022fantastically}\footnote{See Appendix~\ref{sec:diversity} for a rule diversity analysis}. To leverage prediction diversity and promote robustness, bagging is employed by sampling subsets of features or training instances for each run. This is particularly advantageous when training data or feature dimensionality is large. The ensembling strategy confers two principal benefits. Firstly, erroneous rules from some LLM invocations may be offset by accurate rules generated elsewhere, thereby improving the stability of the system. This leverages the self-consistency of LLMs~\cite{wang2022self}, since rules repeatedly inferred across trials tend to be reliable. Secondly, ensemble techniques help mitigate issues caused by prompt length constraints when dealing with large tabular datasets, with bagging reducing prompt size and maintaining performance. If a single rule set overlooks relationships in portions of the data, combining multiple weak learners generated through repeated trials helps alleviate the limitations arising from incomplete coverage in any individual trial.

\section{Experiments}
We assess \textsf{FeatLLM} on a diverse set of tabular datasets and perform multiple analyses to explore the mechanisms underlying our framework. Additional experimental results—including variations on the choice of LLM, qualitative studies, and other aspects—are provided in the Appendix, given space limitations.

\subsection{Performance Evaluation}
\paragraph{Datasets.}
Thirteen datasets, covering binary and multi-class classification scenarios, are selected for experimental evaluation: (1) Adult~\cite{asuncion2007uci}, for inferring if a person earns more than \$50,000 per year; (2) Bank~\cite{moro2014data}, for predicting the likelihood of a customer subscribing to a term deposit; (3) Blood~\cite{yeh2009knowledge}, for identifying donors who will return for further donations; (4) Car~\cite{kadra2021well}, for rating car quality; (5) Communities~\cite{misc_communities_and_crime_183}, for estimating the crime rate within a locality; (6) Credit-g~\cite{kadra2021well}, for distinguishing between favorable and unfavorable credit risks; (7) Diabetes\footnote{\texttt{kaggle.com/datasets/uciml/pima-indians-diabetes-database}}, for diagnosing diabetes in individuals; (8) Heart\footnote{\texttt{kaggle.com/datasets/fedesoriano/heart-failure-prediction}}, for forecasting coronary artery disease occurrence; (9) Myocardial~\cite{misc_myocardial_infarction_complications_579}, for determining chronic heart failure in patients.

Recent additions to the UCI repository as well as synthetic datasets—none of which appeared in the LLMs’ pre-training corpus—are also considered: (10) Cultivars, for estimating the grain yield of a soybean cultivar\footnote{\texttt{archive.ics.uci.edu/dataset/913}}; (11) NHANES, for classifying a person's age group based on their record\footnote{\texttt{archive.ics.uci.edu/dataset/887}}; (12) Sequence-type, which assigns sequences to one of four classes (arithmetic, geometric, Fibonacci, or Collatz); (13) Solution-mix, which evaluates whether the concentration in a four-solution mixture exceeds 0.5. The first two are newly incorporated datasets post-September 2023 in the UCI repository, guaranteeing they were excluded from the pre-training data for the LLM API (gpt-3.5-turbo-0613) applied in this study. The latter two are synthetic, preventing prior exposure in LLM pre-training and facilitating unbiased evaluation.

Further specifics on dataset sizes and complexities are reported in Table~\ref{Tab:basic-desc}.
Every dataset features clearly described attributes. Certain datasets, such as `Communities' and `Myocardial', consist of over a hundred features, posing substantial hurdles related to LLM context window limitations.

\vspace*{-0.12in}
\paragraph{Baselines.}
We benchmark \textsf{FeatLLM} against nine existing methods. The first three are traditional tabular machine learning models: (1) Logistic Regression (LogReg), (2) XGBoost~\cite{chen2016xgboost}, and (3) RandomForest~\cite{ho1995random}. Two additional methods—(4) SCARF~\cite{bahri2022scarf} and (5) STUNT~\cite{nam2023stunt}—assume access to unlabeled data. In practical settings, obtaining substantial unlabeled data may present challenges. (6) TabPFN~\cite{hollmann2022tabpfn}, a recently introduced approach, leverages synthetically generated datasets with varied distributions for model pretraining. The final three baselines adopt LLMs for tabular tasks: (7) In-context learning (In-context)~\cite{wei2022emergent}; (8) TABLET~\cite{slack2023tablet}; and (9) TabLLM~\cite{hegselmann2023tabllm}. In-context learning uses the input prompt to insert a small number of training instances, omitting parameter tuning. TABLET extends this with hints embedded from external classifiers—including rule-sets and prototypes—aiming to improve inference accuracy. TabLLM employs parameter-efficient tuning strategies like IA3~\cite{liu2022few} to adapt the LLM to few-shot samples.

\vspace*{-0.12in}
\paragraph{Implementation details.}
\textsf{FeatLLM} adopts GPT-3.5 as its foundational LLM, though the framework itself is designed to operate independently of the specific LLM chosen (refer to Appendix~\ref{sec:backbones} for comparative results across LLM architectures). The LLM inference process uses a temperature of 0.5 and retains the default API value for top-p at 1. For feature extraction, we configure the ensemble count to 20 and the rule count to 10. Analysis regarding hyper-parameter sensitivity can be found in Figure~\ref{fig:hyperparameter2} and Appendix~\ref{sec:hyper-parameter}. 

For the linear model, we employ the Adam optimizer with a learning rate of 0.01, and carry out training for 200 epochs. Model selection is accomplished via $k$-fold cross-validation, ensuring optimal epoch determination. In the reproduction of baseline experiments, we use GPT-3.5 for in-context learning methods, while T0~\cite{sanh2022multitask} serves as the backbone for TabLLM, consistent with its initial implementation. Notably, TabLLM restricts LLM usage to those offering publicly released checkpoints, which precludes the use of GPT-3.5. Conventional machine learning algorithms, including LogReg, XGBoost, and RandomForest, are tuned by conducting Grid-search in combination with $k$-fold cross-validation; the choice of $k$ (either 2 or 4) ensures that each class is represented in the training folds. For further technical implementation details, please see Appendix~\ref{sec:baselines}.

\vspace*{-0.12in}
\paragraph{Results.}
Tables~\ref{tab:main1} and \ref{tab:main2} present the outcomes of our evaluation across multiple datasets. Each experiment was carried out three times, varying the random selection of training and testing sets; we report both the average and standard deviation of AUC scores. Across all benchmarks, \textsf{FeatLLM} routinely achieves either the highest ranking or finishes as the runner-up when contrasted with alternative baselines. Remarkably, our method reaches similar levels of performance to traditional tabular models using only 4 shots, whereas those require 32 shots (refer to Fig.~\ref{fig:compares}). In some instances, STUNT and TabPFN recorded substantial gains on particular datasets, yet their aggregate improvements over conventional machine learning strategies remained limited. Additionally, Large Language Model-based approaches such as in-context learning, TABLET, and TabLLM were unable to scale inference to tabular datasets featuring extensive numbers of features. Our approach, which leverages both bagging and ensemble strategies, remains effective even on high-dimensional datasets, yielding robust performance. It is worth emphasizing that our concepts of bagging and ensembling can, in principle, be integrated into other LLM-based methods. However, implementing this would necessarily double the inference operations required for each sample, thus significantly increasing computational costs and rendering it unfeasible in practical terms.

\subsection{Analysis \& Discussion}
\label{sec:analysis}
\paragraph{Ablation study.} To elucidate the impact of main architectural elements on performance, we conduct a series of ablation experiments, examining: (1) \textsf{-Tuning}: excluding the weight tuning process from the linear model and instead aggregating binary features to compute the logit; (2) \textsf{-Ensemble}: removing the ensemble step; (3) \textsf{-Description}: not including feature descriptions; and (4) \textsf{-Reasoning}: omitting Step-1 from the reasoning instructions when generating rule sets.

Table~\ref{tab:ablation} reports the average change in AUC following each ablation, computed across all datasets. Omitting any of the key components uniformly leads to a reduction in performance. Notably, the advantage conferred by weight tuning grows as the number of shots increases. This observation suggests that with more available data, rule importance can be estimated more reliably, rendering tuning more impactful. Conversely, the value added by feature descriptions and reasoning instructions is particularly pronounced in low-shot scenarios. This implies that the strategic use of the LLM's prior knowledge is especially beneficial when sample sizes are limited. Finally, our ensemble methodology delivers consistent performance gains regardless of the experimental setting.

\vspace*{-0.12in}
\paragraph{Training \& inference time comparison.} We evaluate and contrast the training and inference durations across various approaches. All experiments are conducted using the Adult dataset, which incorporates a training subset with 16 shots. A single A100 GPU serves as the default hardware configuration, except for TabLLM, which utilizes four A100 GPUs to enable model parallelism. Table~\ref{tab:cost} presents the comparative runtime statistics. Importantly, \textsf{FeatLLM} achieves notably efficient inference times, on par with traditional baselines. Conversely, alternative LLM-based strategies—including In-context learning, TABLET, and TabLLM—are associated with considerably elevated inference times. For training duration, \textsf{FeatLLM} matches the performance of other LLM-based techniques: its requirement of only 30 API calls does not substantially impact the budget and can be easily parallelized.

\vspace*{-0.12in}
\paragraph{Quality of features generated by \textsf{FeatLLM}.}
To assess the value of rule-driven features created by \textsf{FeatLLM} for real-world applications, we carry out an experiment centered on their informativeness. The procedure involves calculating the average AUROC for each generated feature relative to the target labels and then determining the proportion of features with AUROC greater than 0.5 in each dataset. Table~\ref{tab:quality} collates these outcomes. The analysis demonstrates that most of the constructed rules align with the target attribute and deliver valuable information for solving the task (see the ``Before tuning'' column). Furthermore, when the linear model is optimized using the data, features with minimal connection to the dataset (i.e., features whose weights fall below a specified threshold) are systematically eliminated (refer to the ``After tuning'' column). Here, the threshold is fixed at 0.1, based on the distribution of the learned weights.

\vspace*{-0.12in}
\paragraph{Impact of incorporating spurious correlations.} To evaluate the capability of different approaches in removing irrelevant spurious correlations under limited supervision, we performed an empirical study. Specifically, we selected the Heart dataset and arbitrarily merged additional columns sourced from the Adult dataset—these columns have no semantic relevance to the predictive target of the Heart dataset. We tracked model performance as the quantity of extraneous columns increased. For consistency, all methods were given 8 labeled samples, and no unlabeled data was incorporated, thereby omitting methods such as SCARF and STUNT from this evaluation. 

The results, illustrated in Figure~\ref{fig:spurious}, show how model efficacy evolves with the introduction of spurious correlations. Notably, \textsf{FeatLLM} suffered the smallest decline in accuracy compared to standard baselines. This resilience can be attributed to the framework’s requirement that extracted rules must establish clear connections between features and the intended task, guided by prior knowledge—an approach that curtails the influence of irrelevant features and fosters robustness. In practice, we noted that rules originating from noisy features comprised merely 1-2\% of all extracted rules. However, it is also evident that \textsf{FeatLLM} is not immune to learning spurious patterns and misinterpreting causality when columns derived from datasets that bear task-related similarities are appended at random (refer to Appendix~\ref{sec:spurious-appendix}).

\vspace*{-0.12in}
\paragraph{Hyper-parameter analysis}
\textsf{FeatLLM} relies on two principal hyper-parameters: the ensemble count and the rule quantity per class generated for each LLM interaction. We systematically examined their effects on overall performance through a series of experiments. Findings indicate that increasing the number of ensembles results in improved outcomes, though this improvement plateaus after reaching approximately 20 ensembles (refer to Fig.~\ref{fig:appendix-hyperparameter1} in Appendix). With respect to the number of extracted rules per class, no statistically significant variation in performance was observed across settings. However, 10 rules per class was empirically identified as optimal, balancing prompt length and predictive accuracy (see Fig.~\ref{fig:hyperparameter2}). We hypothesize that adding more rules elevates the input dimensionality—while this can provide additional information, it is also prone to inducing overfitting, especially within low-shot contexts.

\section{Conclusion}
This work introduces \textsf{FeatLLM}, a method that sets a new benchmark for few-shot tabular learning using LLMs. Unlike conventional approaches that employ LLMs solely for per-instance predictions, \textsf{FeatLLM} prioritizes generating interpretive criteria reminiscent of feature engineering strategies. By leveraging a rule extraction mechanism together with bagged ensembles, \textsf{FeatLLM} offers low inference demands, operates purely via API interfaces without requiring model retraining, and remains robust to limitations in feature dimensions. This approach delivers strong predictive results and presents an efficient alternative for utilizing LLMs in real-world tabular applications where data is scarce.

\textsf{FeatLLM} is currently tailored for low-shot learning scenarios. In future work, we aim to extend its scope to support feature discovery within datasets possessing more abundant samples, incorporate a broader spectrum of feature types beyond rule-based descriptions, and enhance model interpretability, facilitating its deployment across diverse tasks.

\section{Broader Impact}
The framework detailed here, \textsf{FeatLLM}, harnesses LLMs’ pretrained knowledge to elevate tabular learning performance beyond the reach of conventional supervised learning methodologies. By striving for greater data efficiency, our work addresses overfitting issues inherent in low-sample regimes through the infusion of prior knowledge. \textsf{FeatLLM} offers tangible advantages for practitioners operating in domains where labeling is resource-intensive—such as finance or healthcare—since relevant expertise is embedded within the pretrained LLMs via domain-specific textual data. A critical area for further investigation is assessing how social biases and misinformation encoded within LLM priors may influence outputs; the prior knowledge incorporated into LLMs could impact, or even outweigh, predictions based on available labeled data, raising important considerations for deployment and adoption.

\end{document}